% !TEX root = ../../../main.tex
% Block card: catalogue.snippets -- "Snippets" (Part I, the blocks chapter).
\cardfield{What it does}
Picks the few most typical pieces of a recording, its snippets, and labels every window by the
snippet it most resembles.

\begin{plainterms}[Snippets]
Think of the handful of thumbnail frames that best sum up a film, with a coloured bar under the
timeline showing which thumbnail each scene is closest to. Snippets does that for a recording: cut
it into pieces of one length (an hour, say), find the piece that looks like most of the recording,
then the piece that best covers what the first one does not, and so on. Each window is then
coloured by its nearest snippet. Long runs of one colour suggest a regime of behaviour; a colour
that flickers from window to window suggests none.
\end{plainterms}

\cardfield{When to use it}
When you want a first map of what kinds of behaviour a long recording holds and where each one
dominates, without describing them in advance. Put \blockref{preprocessing.sliding_windows} with no
gap before it (the windows are the pieces), and remove the baseline before that. It is marked
experimental: on the real recordings tried so far, its regimes were mostly no more persistent than
on shuffled copies.

\cardfield{What goes in, what comes out}
In go the windows of a span. Out comes a label per window, 1 to $k$, naming the snippet it most
resembles; $-1$ for a window more than half missing. Each snippet is a real window of the recording,
never an average.

\cardfield{Settings that matter}
\emph{subseq\_frac}, first: how short the sub-pieces are that two pieces are compared by. Smaller
tolerates more shifting and drift inside a piece. \emph{k} is how many snippets; 0 lets the block
choose, and an answer of 1 means it saw no evidence of more than one kind of behaviour. Set \emph{k}
by hand to compare recordings at the same number.

\cardfield{Watch out for}
Only shape counts, so a quiet stretch of plain noise has no shape of its own and is absorbed by
whichever snippet it happens to resemble. Do not trust the automatic \emph{k} on fewer than about 24
pieces (24 hours of hour-long windows). Read the run lengths before calling a label a regime. The
cost grows with the square of the span; spans over a week of 1\,Hz data are refused.
